\subsubsection{The Trace Theorem}\label{subsec:sobolev-trace-theorem}

In this subsection we address the question of whether functions in Sobolev spaces can be restricted to the boundary of the domain, or more generally, to other surfaces. We are also interested in the converse question, namely, how smooth boundary data have to be so that a function in $H^k(\Omega)$ can assume those boundary data. We first consider functions defined on $\mathbb{R}^m$ and their restriction to $\mathbb{R}^{m-1}\times\{0\}$.

\begin{Theorem}
\label{thm:hyperplane-sobolev-trace}
Let $s>1/2$ be real. Then there exists a continuous linear map $T:H^s(\mathbb{R}^m)\to H^{s-1/2}(\mathbb{R}^{m-1})$, called the trace operator, with the property that for any $\phi\in\mathcal{D}(\mathbb{R}^m)$, we have
\begin{equation}\label{eq:hyperplane-trace-test-function-restriction}
T\phi(x_1,\ldots,x_{m-1})=\phi(x_1,x_2,\ldots,x_{m-1},0).
\end{equation}

\end{Theorem}

It can be shown that for $s\leq1/2$ the result is false; in fact, functions which vanish in a neighborhood of $x_m=0$ are then dense in $H^s(\mathbb{R}^m)$.

\begingroup
\postdisplaypenalty=10000
\renewcommand{\endProof}{\par\nobreak\hfill$\blacksquare$\par\vspace{0.1cm}}
\begin{Proof}
Let $\phi\in\mathcal{D}(\mathbb{R}^m)$ and let $g(\bm{x}^{\prime})=\phi(\bm{x}^{\prime},0)$. Let $\widetilde\phi$ denote the Fourier transform of $\phi$ with respect to the $m$th variable only, i.e.,
\begin{equation}\label{eq:trace-proof-partial-normal-fourier-transform}
\widetilde\phi(\bm{x}^{\prime},\xi_m)=\dfrac{1}{\sqrt{2\pi}}\int_{-\infty}^{\infty}\phi(\bm{x}^{\prime},x_m)e^{-ix_m\xi_m}\,\mathrm{d}x_m,
\end{equation}
let $\hat\phi$ and $\hat g$ denote the Fourier transforms of $\phi$ and $g$ in $\mathbb{R}^m$ and, respectively, $\mathbb{R}^{m-1}$. The Fourier inversion formula yields
\begin{equation}\label{eq:trace-proof-normal-fourier-inversion}
g(\bm{x}^{\prime})=\phi(\bm{x}^{\prime},0)=\dfrac{1}{\sqrt{2\pi}}\int_{-\infty}^{\infty}\widetilde\phi(\bm{x}^{\prime},\xi_m)\,\mathrm{d}\xi_m.
\end{equation}
Applying the Fourier transform, we find
\begin{equation}\label{eq:trace-proof-tangential-fourier-integral}
\hat g(\bm{\xi}^{\prime})=\dfrac{1}{\sqrt{2\pi}}\int_{-\infty}^{\infty}\hat\phi(\bm{\xi})\,\mathrm{d}\xi_m.
\end{equation}
We now estimate
\begin{equation}\label{eq:trace-proof-weighted-fourier-cauchy-estimate}
\begin{aligned}
\|g\|_{s-1/2,2}^2&\leq C_1\int_{\mathbb{R}^{m-1}}|\hat g(\bm{\xi}^{\prime})|^2(1+|\bm{\xi}^{\prime}|^2)^{s-1/2}\,\mathrm{d}\bm{\xi}^{\prime}\\
&=\dfrac{C_1}{2\pi}\int_{\mathbb{R}^{m-1}}\left|\int_{-\infty}^{\infty}\hat\phi(\bm{\xi})\,\mathrm{d}\xi_m\right|^2(1+|\bm{\xi}^{\prime}|^2)^{s-1/2}\,\mathrm{d}\bm{\xi}^{\prime}\\
&\leq C_2\int_{\mathbb{R}^{m-1}}(1+|\bm{\xi}^{\prime}|^2)^{s-1/2}\int_{-\infty}^{\infty}|\hat\phi(\bm{\xi})|^2(1+|\bm{\xi}|^2)^s\,\mathrm{d}\xi_m\\
&\qquad\times\int_{-\infty}^{\infty}(1+|\bm{\xi}|^2)^{-s}\,\mathrm{d}\xi_m\,\mathrm{d}\bm{\xi}^{\prime}.
\end{aligned}
\end{equation}
If $s>1/2$, we have
\begin{equation}\label{eq:trace-proof-integrated-normal-fourier-weight}
\begin{aligned}
\int_{-\infty}^{\infty}(1+|\bm{\xi}|^2)^{-s}\,\mathrm{d}\xi_m
&=\int_{-\infty}^{\infty}(1+|\bm{\xi}^{\prime}|^2+\xi_m^2)^{-s}\,\mathrm{d}\xi_m\\
&=(1+|\bm{\xi}^{\prime}|^2)^{-s+1/2}\int_{-\infty}^{\infty}(1+y^2)^{-s}\,\mathrm{d}y.
\end{aligned}
\end{equation}
By inserting into \eqref{eq:trace-proof-weighted-fourier-cauchy-estimate}, we find that $\|T\phi\|_{s-1/2,2}\leq C\|\phi\|_{s,2}$ for every $\phi\in\mathcal{D}(\mathbb{R}^m)$. Hence $T$ can be extended by continuity to all of $H^s(\mathbb{R}^m)$.
\end{Proof}
\endgroup

\begin{Lemma}
\label{lem:continuous-function-trace-and-support}
If $u\in H^s(\mathbb{R}^m)\cap C(\mathbb{R}^m)$, $s>1/2$, then $Tu(\bm{x}^{\prime})=u(\bm{x}^{\prime},0)$. Moreover $\operatorname{supp}Tu\subset\operatorname{supp}u\cap(\mathbb{R}^{m-1}\times\{0\})$ for every $u\in H^s(\mathbb{R}^m)$.
\end{Lemma}

We leave the proof as an exercise.

A natural question is now which elements of $H^{s-1/2}(\mathbb{R}^{m-1})$ can be obtained as restrictions, or ``traces'' of functions in $H^s(\mathbb{R}^m)$. The answer is that all elements of $H^{s-1/2}(\mathbb{R}^{m-1})$ are obtained in this way.
